\section{A new 215-state marker allocation}
\label{sec:marker215}

Let $K$ denote the repository's primary accepted K215 table with SHA-256
\texttt{9683d6ed9628e6ce407c315b08f5da3d063da94e92706f275ca0f33e8272e4b2}.
The construction below changes six transition rows, retains exactly the same
215 working labels and 430 binary transitions, and leaves the source program
unchanged.  The new table $M$ has SHA-256
\texttt{e5907f08e0752eada0e25bbf33680341e80e5fe94b52b409aea1ed039edda252}.
Start, dispatcher and halt retain native labels $0,84,254$ respectively.
The dense version starts at $0$, uses working labels $0,\ldots,214$, and halts
at label $215$.

The idea is to let the tape cell immediately preceding the head retain part
of the logical control state.  Logical roles 52 and 176 become one role at
physical state 52.  The old physical label 176 is reused as the marker reader.
Unlike an unconditional state merge, the proof must account for the exceptional
incoming edge from state 156.  That edge is harmless precisely because it
forces the current cell to the branch on which both roles already agree.

\subsection{The six-row construction}

A transition is written $w D q$, with write symbol $w$, direction $D$, and
next control state $q$.  The changed rows are
\[
\begin{array}{c|c|c|c}
 q&\text{read}&K&M\\\hline
 52&0&0R154&0L176\\
 156&1&1R176&1R52\\
 165&0&0R176&0R52\\
 176&0&0R168&0R185\\
 176&1&1R150&1R129\\
 193&0&0R176&0R52
\end{array}
\]
All unlisted rows are identical.  In particular,
\[
 K(52,1)=K(176,1)=M(52,1)=1R150,
 \qquad M(185,0)=0R168,\quad M(129,0)=0R154.
\]
The complete incoming-edge census in $K$ is
\[
\begin{array}{c|l}
\text{destination}&\text{incoming transitions}\\\hline
52&190/1=1R52,\quad 193/1=1R52\\
176&156/1=1R176,\quad 165/0=0R176,\quad193/0=0R176\\
156&189/1=1L156.
\end{array}
\]
These are exhaustive statements about the literal table, not observations
from a finite blank run.

\subsection{A reachable-branch invariant}

\paragraph{Lemma.}
For every run of $K$ starting in control state 0, on an arbitrary initial
two-way binary tape and at an arbitrary initial head position:
\begin{enumerate}
\item every entry into state 52 has a one immediately to the left;
\item every entry into state 176 which scans zero has a zero immediately to
      the left.
\end{enumerate}

\paragraph{Proof.}
The first assertion follows from the two incoming transitions to 52, both
of which write one and move right.  The incoming transitions $165/0$ and
$193/0$ similarly write zero and move right before entering 176.
The remaining incoming edge is $156/1$.  Since the run starts at 0, every
entry into 156 has a predecessor; its unique predecessor is $189/1$, which
writes one and moves left.  If 156 then follows its one branch, it moves right
back onto the very cell just written by 189, so the current cell at entry to
176 is one.  This exceptional edge can therefore never produce the zero
branch under discussion.  No assumption about uninspected tape cells is
used.\hfill$\square$

\subsection{Exact all-tape simulation from the designated start}

\paragraph{Theorem.}
For every initial binary tape $\tau$ and every initial head position $h$, the
runs of $K$ and $M$ starting at $(0,h,\tau)$ halt simultaneously in the sense
of equihalting.  More precisely, there is a strictly increasing sequence of
simulation times $t_n$ with $t_0=0$ such that the $n$th configuration of $K$
and the $t_n$th configuration of $M$ have exactly the same whole tape and head,
and their controls are related by
\[
 \phi(q)=\begin{cases}52,&q=176,\\q,&q\ne176.\end{cases}
\]
Each increment $t_{n+1}-t_n$ is one or three.  It is three exactly when $K$
executes row $52/0$ or $176/0$.

\paragraph{Proof.}
For every $q\notin\{52,176\}$, the new transition equals the old transition
followed by $\phi$ on its target.  Thus these source steps transfer in one
step, without changing the comparison tape or head.  The one branches of
52 and 176 agree and also transfer in one step.

It remains to consider a zero branch.  If the old state is 52, the lemma
provides predecessor marker one.  The new machine executes
\[
 52 \xrightarrow{0L}176
 \xrightarrow{1R}129
 \xrightarrow{0R}154.
\]
If the old state is 176, the predecessor marker is zero and the new machine
executes
\[
 52 \xrightarrow{0L}176
 \xrightarrow{0R}185
 \xrightarrow{0R}168.
\]
Each block reads only the original current cell and its predecessor,
preserves both symbols, and finishes one cell to the right.  It therefore
has exactly the whole-tape effect of the corresponding old transition.
Induction on the old run length proves the displayed simulation relation.
Both inserted blocks have finite positive length and contain no halt.
Consequently a finite halt transfers in both directions; an infinite old
run produces infinitely many new steps; and the new run cannot remain
forever within one inserted block.\hfill$\square$

If a finite old execution of length $n$ has $v(n)$ visits to the two specified
zero rows, its corresponding new execution has the exact length
\[
 t_n=n+2v(n),\qquad n\le t_n\le3n.
\]
The theorem deliberately quantifies arbitrary tapes only at the designated
start control.  It does not claim equivalence from arbitrary controls such
as 156 or 176, where the predecessor invariant need not hold.

\paragraph{Consistency calibration.}
The preceding theorem is a finite transition argument followed by elementary
induction, formalizable in ZFC.  Combining it with the repository's accepted
ordinary calibration of $K$ gives
\[
 \mathrm{ZFC}\vdash
 \bigl(\neg\operatorname{Halt}(M,\mathrm{blank})\bigr)
 \;\longleftrightarrow\;\operatorname{Con}(\mathrm{ZFC}).
\]
This transfers the accepted calibration, including its underlying numerical
calculus and compiler.  It is not an assertion of unconditional nonhalting
or a claim that a new complete Lean/Rocq endpoint has been kernel checked.

\subsection{Fresh independent checks and measured effect}

The complete inherited symbolic-runtime checker was rerun on the new literal
bytes and unchanged source, rather than relying solely on the transfer theorem.
It proved all 6,176 expected arbitrary-natural primitive cases using
1,701,542 symbolic rules.  All 430 transition rows are accounted for, every
affine duration is positive, and the minimum constant duration is six.
The represented domain includes arbitrary natural primary-register values,
an arbitrary natural gap, and every nonempty finite spare list, with the
entire zero exterior included.  These are quantified mathematical checks,
not bounded register-value tests; they remain outside a proof-assistant kernel.

The whole blank startup takes 95,714 transitions and ends at native dispatcher
84, head $-8$, header 57, gap one, logical registers
$(4,0,0,0,0,0,0,0,0)$, and forty unit spares.  The complete terminal tape has
95 one-cells.  An independently written literal checker replays the stored
trace on both the native and dense tables, constructs the expected whole tape
from the frame parameters, checks the reversible dense-state map, and verifies
the finite incoming-edge and local-block hypotheses above.

Relative to the accepted primary case records, the durations of 5,632 primitive
cases are unchanged, 128 cases gain two steps, and 416 cases gain six steps.
The coefficients of every unbounded affine variable and all net head motions
are unchanged.  Thus 544 quantified primitive cases use a changed execution,
even though the finite blank startup happens not to encounter the folded zero
branches.

The new table has one more left-moving transition than the primary table,
which already rules out a pure state renaming.  A separate exact rooted
deterministic transition-graph comparison, also allowing global tape
reflection, found no match among 77 distinct 215-state tables recovered from
twelve incoming research archives and the accepted candidate directory.
This is a comparison against that explicit audited corpus, not a claim of
exhaustive novelty among all possible machines.

\subsection{Why the construction is useful for the 214-state target}

This allocation is cost-neutral: one physical role is freed by sharing the
two logical roles, and that same label is spent on a new marker selector.
It is therefore a genuine 215-state alternative rather than a 214-state
claim.  A later source reassociation that makes an existing state implement
$0R185/1R129$ at all required contexts would remove the need for the extra
selector label.  More generally, the lemma permits an incoming marker anomaly
when that anomaly forces the common branch.  Requiring every incoming edge
to carry the same marker would miss precisely this opportunity.
